% !TEX root = ../main.tex
\section{Implementation}
\label{sec:implementation}

The implementation is written from scratch in PyTorch; no reference code for
FedNPG-ADMM is publicly available. The algorithm proper (\texttt{agent.py},
\texttt{server.py}) stays free of orchestration concerns, and every
agent-to-server transmission is routed through a single \texttt{Channel}
abstraction that meters bytes. That chokepoint is what makes the communication
comparison in Section~\ref{sec:baselines} meaningful: all three algorithms
count bytes through the same code, and the quantization study of
Section~\ref{sec:quantization} attaches there without any algorithm importing
it.

\subsection{What comes from the paper and what does not}

Table~\ref{tab:table4} reproduces the paper's hyperparameters. Two values in
our pipeline are \emph{not} transcriptions. First, $\eta$ in line~12 of
Algorithm~\ref{alg:fednpgadmm} appears nowhere in Table~4 of the original. The
paper derives it analytically from smoothness constants it never reports
numerically, so we set $\eta = 1$, taking the ADMM step unscaled. Second,
Table~4's ``Learning rate'' column ($3\times10^{-4}$, $3\times10^{-4}$,
$10^{-5}$) is used here as the \emph{critic's} Adam learning rate, not an
actor step size. Those are characteristic Adam values, and the paper defines
the actor's step analytically, which we read as evidence that the column
refers to the critic. This is an interpretation of an ambiguous table.

\subsection{Practical settings}

The paper's own setup section states that it follows ``the practical settings
in TRPO with line search and in stable-baselines3 with generalized advantage
estimation (0.95) and the Adam optimizer''. Algorithm~\ref{alg:fednpgadmm}'s
pseudocode contains none of these, so we adopt the standard versions and list
them in Appendix~\ref{app:hyperparameters}: running observation and reward
normalization, advantage normalization, orthogonal initialization, Generalized Advantage
Estimation (GAE) with $\lambda = 0.95$, minibatch critic fitting, and a backtracking line search on
the true KL and the importance-sampled surrogate. The line search matters. The
closed-form step of line~12 controls only the second-order approximation of
the KL, so without a safeguard it can overshoot the intended trust region.

Two further design choices are ours because the paper is silent. The critic is
fit on data \emph{pooled} across agents; the paper describes only the actor's
update and never says how a federated critic is trained. And when $(\sum_i
\bg_i)^{\top}\by \le 0$ the step size in line~12 is undefined, in which case
we skip the update for that iteration and log how often it happens.

\subsection{Reporting conventions}

We report the paper's own metric: the average trajectory return across the
batch collected at each iteration, plotted against iterations. Two measurement
issues are worth stating because they materially changed our numbers.

First, this metric averages only the episodes that \emph{complete} within one
iteration's timestep budget. On Hopper at $N=1$ that is roughly $1.2$ episodes
per point, so each plotted value is close to a single episode, which makes the
curve appear unstable even when the underlying policy is not. We therefore
also record a trailing 100-episode mean and, separately, a deterministic
evaluation (mean action, no exploration noise, on a held-out environment), the
quantity published benchmarks such as \texttt{rl-baselines3-zoo} report. The
deterministic number runs roughly $1.2\times$ higher than the stochastic one,
so the two are not interchangeable.

Second, no published MuJoCo learning curve plots this metric raw.
\citet{spinningup} smooth the identical quantity over an 11-epoch window and
average 10 seeds, and \texttt{rl-baselines3-zoo} uses a 100-episode rolling
window. Where we show curves we state the smoothing in the caption, following
the recommendation of \citet{huang2024openrlbenchmark}.
